%%%%%
%
\section{Calcul des flux de chaleur sensible par m�thode a�rodynamique}

\itemstyle{Op\'erateur~:} {\sf Flux\_Aero}
\index{Flux\_Aero}

%%%%%OBJET%%%%
\itemstyle{Objet~:}
 Cet op�rateur impl�mente la m�thode dite a�rodynamique,
qui utilise les relations liant, d'une part $u_*$ et
$du \over dz$, et d'autre part $\theta_*$ et $d\theta \over dz$,
(profils universels de la th�orie de la similitude de Monin-Obukhov)
pour d�duire le flux de chaleur sensible ($-\rho c_p u_* \theta_*$)
� partir de mesures des gradients verticaux de temp�rature et de vent.

La hauteur s�parant les 2 niveaux de mesure est
consid�r�e fixe (1.5 m), mais la hauteur au-dessus du sol, et 
la hauteur de la v�g�tation (pour estimer la hauteur de
d�placement) sont des variables d'entr�e.

En fin de calcul, le flux de chaleur latente (LE) est
estim� en supposant le bilan �nerg�tique ferm�.

%%%%%ARGUMENTS%%%%
\bigskip
%%%%%ENTREES%%%%
\begin{tabarg}
\multicolumn{3}{l}{\titretab {Entr\'ees}}\\
\hline
$dv$ & Vecteur, {\tt float} & ffhaut-ffbas (m/s)\\
$ff_{hau}$ & Vecteur, {\tt float} & vitesse du vent superieur (m/s)\\
$dt_{air}$ & Vecteur, {\tt float} & $t_{hau}-t_{bas}$ (celsius)\\
$t$ & Vecteur, {\tt float} & temp�rature sous abri (celsius)\\
$rn$ & Vecteur, {\tt float} & rayonnement net (W/m2)\\
$phi0$ & Vecteur, {\tt float} & flux dans le sol (W/m2)\\
$hveg$ & Vecteur, {\tt float} & hauteur de v�g�tation (m)\\
$hbas$ & Vecteur, {\tt float} & hauteur niveau inf. des mesures (m)\\
$p$ & Vecteur, {\tt float} & pression (hPa)\\
$humid$ & Vecteur, {\tt float} & humidit� (\%)\\
\end{tabarg}

\medskip
%%%%%SORTIES%%%%
\begin{tabarg}
\multicolumn{3}{l}{\titretab {Sorties}}\\
\hline
$hs$ & Vecteur, {\tt float} & flux de chaleur sensible (W/m2)\\
$le$ & Vecteur, {\tt float} & flux de chaleur latente (W/m2)\\
$z_0$ & Vecteur, {\tt float} & rugosite (m)\\
$vlmon$ & Vecteur, {\tt float} & longueur de Monin Obukhov (m)\\
$ustar$ & Vecteur, {\tt float} & vitesse de frottement (m/s)\\
$tstar$ & Vecteur, {\tt float} & temp�rature de frottement (Celsius)\\
\end{tabarg}


%%%%%FORMULE%%%%
\itemstyle{Formule}

%%Flux_Aero

selon 3 cas de stabilite : "convection libre" / neutre / autre.
Resolution iterative jusqu'� ce que deux valeurs successives de la
longueur de Monin Obukhov diff�rent de moins de 5 % (seuil choisi).
La convergence du processus it�ratif n'est assur� que pour des conditions
instables.

Cas de la convection libre :
\begin{eqnarray*}
ustar = 0.\\
tstar = -999.\\
vlmon = 0.\\
vhs = 0.\\
\end{eqnarray*}

cas neutre :

\begin{eqnarray*}
ustar = frac{K.dv}{log(zu_2/zu_1)}\\
tstar = 0.\\
vlmon = 9999.0.sign(1.,dt_air)\\
vhs = 0.\\
\end{eqnarray*}

autres cas (stratification neutre, soit vlmon infinie) :

\begin{eqnarray*}
vlmon = 1000000.\\
ustar = frac{K . dv}{log(zu_2/zu_1) - psi_m(zu_2/vlm_1) + psi_m(zu_1/vlm_1)}\\
tstar = frac{K . dt_air}{log(zu_2/zu_1) - psi_h(zu_2/vlm_1) + psi_h(zu_1/vlm_1)}\\
vlmon = frac{tk . ustar^2}{karman . g . tstar}\\
\end{eqnarray*}\\
 
Calcul du flux de chaleur sensible vhs :

Cas instable (avec convergence)

vhs = -roo . cp . ustar . tstar

rho : calcul de la masse volumique d'une particule d'air humide
definie par : pression , temperature , humidite
rho     : reel   : masse en kg/m**3
pres   : reel   : pression em mb
temp   : reel   : temperature en celsius
humi   : reel   : humidite en %

fonction tvir\_4m(pres,temp,humi)
calcul de la temp\'erature virtuelle d'une particule
donnee definie par : pression , temperature , humidite
tvir\_4m : reel  : temp\'erature en degr\'es Celsius

%%%
%% Reste � d�crire l'it�ration. En g�n�ral, c'est
%% sous la forme :
%% soit F(x) dont on chreche les z�ros
%% 1) estimation de X0
%% 2) Calcul de F(Xi)
%% 3) Calcul de F'(Xi)
%% 4) Calcul de Xi+1 par une m�thode
%% (Si Newton, xi+1 = Xi - F(xi)/F'(xi) )
%%  5) si pas convergence, bouclage en 2
%%
Cas stable (sans convergence)

$vhs = {-karman^2 . roo . cp . dv . dt_air}{\log(zu_2/zu_1)^2 / 10}$




avec~:\\
$t_0 = 273.15 $ \\
$karman = 0.4 $ \\
$cp = 1005.0 $ \\
$g = 9.81 $ \\

Dyer 74


%%%%COMMENTAIRE%%%%
\itemstyle{Commentaire~:}

Les fonctions $rho()$ et $tvir\_4m()$ sont d\'ecrites en annexe\ref{annexe-thermo}.

%%%%%AUTEUR%%%%
\itemstyle{Auteur~:}
CNRM/GMEI/TRAMM

%%%%%REFERENCE%%%%
\itemstyle{R\'ef\'erences~:}
Programme Fortran 77 de Dominique Puech et note de centre CNRM/GMEI nr 6 \cite{GMEI-6}.
